\subsection{THE JORDAN-HÖLDER THEOREM}

DEFINITION 9. A composition series of a group with operators G is a finite sequence \((\mathbf{G}_i)_{0\leqslant i\leqslant n}\) of stable subgroups of G, with \(\mathbf{G}_0 = \mathbf{G}\) and \(\mathbf{G}_n = \{e\}\) and such that \(\mathbf{G}_{i + 1}\) is a normal subgroup of \(\mathbf{G}_i\) for \(0\leqslant i\leqslant n - 1\) . The quotients \(\mathrm{G_i / G_{i + 1}}\) are called the quotients of the series. A composition series \(\Sigma^{\prime}\) is said to be finer than a composition series \(\Sigma\) if \(\Sigma\) is a series taken from \(\Sigma^{\prime}\) .

If \((\mathbf{G}_i)_{0\leqslant i\leqslant n}\) and \((\mathbf{H}_j)_{0\leqslant j\leqslant m}\) are respectively composition series of two groups with operators \(\mathbf{G}\) and \(\mathbf{H}\) , they are said to be equivalent if \(m = n\) and there exists a permutation \(\phi\) of the interval \([0,n - 1]\) of \(\mathbf{N}\) , such that the groups with operators \(\mathbf{G}_i / \mathbf{G}_{i + 1}\) and \(\mathrm{H}_{\phi (i)} / \mathrm{H}_{\phi (i) + 1}\) are isomorphic for all \(i\) .

Note that in general a series taken from a composition series \((\mathbf{G}_{i})\) is not a composition series, since for \(j > i + 1\) , \(G_{j}\) is not in general a normal subgroup of \(G_{i}\) .

THEOREM 5 (Schreier). Given two composition series \(\Sigma_{1}, \Sigma_{2}\) of a group with operators G, there exist two equivalent composition series \(\Sigma_{1}^{\prime}, \Sigma_{2}^{\prime}\) , finer respectively than \(\Sigma_{1}\) and \(\Sigma_{2}\) .

Let \(\Sigma_{1} = (\mathrm{H}_{i})_{0\leqslant i\leqslant n}\) and \(\Sigma_{2} = (\mathrm{K}_{j})_{0\leqslant j\leqslant p}\) be the two given composition series with respectively \(n + 1\) and \(p + 1\) terms; we shall see that the composition series \(\Sigma_1^\prime\) can be formed by inserting \(p - 1\) subgroups \(\mathbf{H}_{i,j}^{\prime}\) \((1 \leqslant j \leqslant p - 1)\) between \(H_{i}\) and \(H_{i+1}\) for \(0 \leqslant i \leqslant n - 1\) and the series \(\Sigma_{2}'\) by inserting \(n - 1\) subgroups \(K_{j,i}'\) ( \(1 \leqslant i \leqslant n - 1\) ) between \(K_{j}\) and \(K_{j+1}\) for \(0 \leqslant j \leqslant p - 1\) ; thus two series of \(pn + 1\) stable subgroups of G will be obtained; by choosing suitably the inserted stable subgroups, we shall show that these series are equivalent composition series.

To this end note that \(\mathbf{H}_i \cap \mathbf{K}_j\) is a stable subgroup of \(\mathbf{H}_i\) and of \(\mathbf{K}_j\) and hence (Theorem 4) \(\mathbf{H}_{i+1} \cdot (\mathbf{H}_i \cap \mathbf{K}_j)\) is a stable subgroup of \(\mathbf{H}_i\) containing \(\mathbf{H}_{i+1}\) and \(\mathbf{K}_{j+1} \cdot (\mathbf{H}_i \cap \mathbf{K}_j)\) is a stable subgroup of \(\mathbf{K}_j\) containing \(\mathbf{K}_{j+1}\) . If we write \(\mathbf{H}_{i,j}' = \mathbf{H}_{i+1} \cdot (\mathbf{H}_i \cap \mathbf{K}_j)\) and \(\mathbf{K}_{j,i}' = \mathbf{K}_{j+1} \cdot (\mathbf{H}_i \cap \mathbf{K}_j)\) , \(\mathbf{H}_{i,j+1}'\) is a stable subgroup of \(\mathbf{H}_{i,j}'\) ( \(0 \leqslant j \leqslant p-1\) ) and \(\mathbf{K}_{j,i+1}'\) is a stable subgroup of \(\mathbf{K}_{j,i}'\) ( \(0 \leqslant i \leqslant n-1\) ). Moreover \(\mathbf{H}_{i,0}' = \mathbf{H}_i\) , \(\mathbf{H}_{i,p}' = \mathbf{H}_{i+1}\) , \(\mathbf{K}_{j,0}' = \mathbf{K}_j\) and \(\mathbf{K}_{j,p}' = \mathbf{K}_{j+1}\) . To show the theorem, it suffices to show that \(\mathbf{H}_{i,j+1}'\) (resp. \(\mathbf{K}_{j,i+1}'\) ) is a normal stable subgroup of \(\mathbf{H}_{i,j}'\) (resp. \(\mathbf{K}_{j,i}'\) ) and that the quotient groups \(\mathbf{H}_{i,j}'/\mathbf{H}_{i,j+1}'\) and \(\mathbf{K}_{j,i}'/\mathbf{K}_{j,i+1}'\) are isomorphic ( \(0 \leqslant i \leqslant n-1\) , \(0 \leqslant j \leqslant p-1\) ). This follows from the following lemma by taking \(\mathbf{H} = \mathbf{H}_i\) , \(\mathbf{H}' = \mathbf{H}_{i+1}\) , \(\mathbf{K} = \mathbf{K}_j\) , \(\mathbf{K}' = \mathbf{K}_{j+1}\) .

Lemma 1 (Zassenhaus). Let H and K be two stable subgroups of a group with operators G and H' and K' normal stable subgroups of H and K respectively; then H'.(H ∩ K') is a normal stable subgroup of H'.(H ∩ K), K'.(K ∩ H') is a normal stable subgroup of K'.(K ∩ H) and the quotient groups with operators (H'.(H ∩ K))/(H'.(H ∩ K')) and (K'.(K ∩ H))/(K'.(K ∩ H')) are isomorphic.

By Theorem 4, \(\mathbf{H}' \cap \mathbf{K} = \mathbf{H}' \cap (\mathbf{H} \cap \mathbf{K})\) is a normal stable subgroup of \(\mathbf{H} \cap \mathbf{K}\) ; similarly \(\mathbf{K}' \cap \mathbf{H}\) is a normal stable subgroup of \(\mathbf{K} \cap \mathbf{H}\) ; hence (no. 6, Corollary 2) \((\mathbf{H}' \cap \mathbf{K})(\mathbf{K}' \cap \mathbf{H})\) is a normal stable subgroup of \(\mathbf{H} \cap \mathbf{K}\) . By Theorem 4 applied to the group \(\mathbf{H}\) ,

\[
\mathbf {H} ^ {\prime}. (\mathbf {H} ^ {\prime} \cap \mathbf {K}). (\mathbf {K} ^ {\prime} \cap \mathbf {H}) = \mathbf {H} ^ {\prime}. (\mathbf {H} \cap \mathbf {K} ^ {\prime})
\]

is a normal stable subgroup of \(\mathbf{H}'\) . (H ∩ K) and the quotient group

\[
(\mathbf {H} ^ {\prime}. (\mathbf {H} \cap \mathbf {K})) / (\mathbf {H} ^ {\prime}. (\mathbf {H} \cap \mathbf {K} ^ {\prime}))
\]

is isomorphic to

\[
(\mathbf {H} \cap \mathbf {K}) / ((\mathbf {H} ^ {\prime} \cap \mathbf {K}). (\mathbf {K} ^ {\prime} \cap \mathbf {H})).
\]

In the last quotient, H and H' on the one hand and K and K' on the other appear symmetrically; permuting them the stated result is obtained.

DEFINITION 10. A Jordan-Hölder series of a group with operators G is a strictly decreasing decomposition series \(\Sigma\) such that there exists no strictly decreasing decomposition series distinct from \(\Sigma\) and finer than \(\Sigma\) .

PROPOSITION 9. For a decomposition series of G to be a Jordan-Hölder series it is necessary and sufficient that all the quotients of the series be simple.

A decomposition series is strictly decreasing if and only if none of its successive quotients is reduced to the identity element. If a strictly decreasing composition series \(\Sigma\) is not a Jordan-Hölder series, there exists a strictly decreasing composition series \(\Sigma'\) which is finer than \(\Sigma\) and distinct from \(\Sigma\) . There are therefore two consecutive terms \(G_i\) , \(G_{i+1}\) of \(\Sigma\) which are not consecutive in \(\Sigma'\) ; let H be the first term which follows \(G_i\) in \(\Sigma'\) ; H is a normal stable subgroup of \(G_i\) , containing \(G_{i+1}\) and distinct from the latter; hence H/ \(G_{i+1}\) is a normal stable subgroup of \(G_i/G_{i+1}\) , distinct from the latter and from the identity element; therefore \(G_i/G_{i+1}\) is not simple. Conversely, if \(\Sigma\) is a strictly decreasing composition series one of whose quotients \(G_i/G_{i+1}\) is not simple, this quotient contains a normal stable subgroup other than itself and \(\{e\}\) , whose inverse image in \(G_i\) is a normal stable subgroup H of \(G_i\) , distinct from \(G_i\) and \(G_{i+1}\) (Theorem 4); it suffices to insert H between \(G_i\) and \(G_{i+1}\) to obtain a strictly decreasing composition series distinct from \(\Sigma\) and finer than \(\Sigma\) .

THEOREM 6 (Jordan-Hölder). Two Jordan-Hölder series of a group with operators are equivalent.

Let \(\Sigma_{1}, \Sigma_{2}\) be two Jordan-Hölder series of a group with operators G; by applying Theorem 5 two equivalent composition series \(\Sigma_{1}^{\prime}, \Sigma_{2}^{\prime}\) are obtained which are respectively finer than \(\Sigma_{1}\) and \(\Sigma_{2}\) ; since the latter are Jordan-Hölder series, \(\Sigma_{1}^{\prime}\) is identical with \(\Sigma_{1}\) or is derived from it by repeating certain terms; the series of quotients of \(\Sigma_{1}^{\prime}\) is derived from that for \(\Sigma_{1}\) by inserting a number of terms isomorphic to the group \(\{e\}\) ; since \(\Sigma_{1}\) is strictly decreasing, the series of quotients of \(\Sigma_{1}\) is derived from that of \(\Sigma_{1}^{\prime}\) by suppressing in the latter all the terms isomorphic to \(\{e\}\) . Similarly for \(\Sigma_{2}\) and \(\Sigma_{2}^{\prime}\) . As the series of quotients of \(\Sigma_{1}^{\prime}\) and \(\Sigma_{2}^{\prime}\) differ (up to isomorphism) only in the order of the terms, the same is true of \(\Sigma_{1}\) and \(\Sigma_{2}\) ; the theorem is proved.

COROLLARY. Let G be a group with operators in which there exists a Jordan-Hölder series. If \(\Sigma\) is any strictly decreasing composition series of G, there exists a Jordan-Hölder series finer than \(\Sigma\) .

Let \(\Sigma_0\) be a Jordan-Hölder series of \(G\) ; by Theorem 5, there exist two equivalent composition series, \(\Sigma'\) and \(\Sigma_0'\) respectively finer than \(\Sigma\) and \(\Sigma_0\) ; the argument of Theorem 6 shows that, by suppressing from \(\Sigma'\) the repetitions, a sequence \(\Sigma''\) is obtained which is equivalent to \(\Sigma_0\) and hence a Jordan-Hölder series, since all its quotients are simple (Proposition 9). As \(\Sigma\) is strictly decreasing, \(\Sigma''\) is finer than \(\Sigma\) , whence the corollary.

Remark. A group with operators does not always possess a Jordan-Hölder series; an example is given by the additive group \(\mathbf{Z}\) of rational integers: the sequence \((2^{n}.\mathbf{Z})_{n\geqslant 0}\) is a strictly decreasing infinite sequence of (normal) subgroups of \(\mathbf{Z}\) ; for all \(p\) , the first \(p\) terms of this sequence form with the group \(\{0\}\) a strictly decreasing composition series; if there existed a Jordan-Hölder series for Z, it would have at least \(p + 1\) terms, by the Corollary to Theorem 6; absurd, since p is arbitrary.

On the other hand, there exists a Jordan-Hölder series in every finite group with operators G: if \(G \neq \{e\}\) , among the normal stable subgroups of G distinct from G, let \(H_{1}\) be a maximal subgroup; similarly define \(H_{n+1}\) by induction as a maximal element in the set of normal subgroups of \(H_{n}\) distinct from \(H_{n}\) , when \(H_{n} \neq \{e\}\) ; the sequence of orders of the \(H_{n}\) is strictly decreasing, hence there exists n such that \(H_{n} = \{e\}\) and the sequence consisting of G and the \(H_{i} (1 \leqslant i \leqslant n)\) is by its formation a Jordan-Hölder series.

DEFINITION 11. Let G be a group with operators; the length of G is the upper bound of the integers n such that there exists a strictly decreasing composition series of G (G \(_{i}\) ) \(_{0 \leq i \leq n}\) .

If G admits a Jordan-Hölder series, the length of G is the number of successive quotients of this series, as follows from the Corollary to Theorem 6. If G does not admit a Jordan-Hölder series, its length is infinite; by Proposition 9, for every strictly decreasing series of G there exists a strictly finer strictly decreasing series. The group consisting of the identity element is the only group with operators of length zero. A group with operators is simple if and only if it is of length 1.

Let G be a group with operators, H a normal stable subgroup of G, K the quotient G/H and \(\pi: G \rightarrow K\) the canonical homomorphism. Let

\[
\Sigma^ {\prime} = \left(\mathrm{H} _ {i}\right) _ {0 \leqslant i \leqslant n}
\]

be a decomposition series of H and \(\Sigma'' = (\mathbf{K}_j)_{0 \leqslant j \leqslant p}\) be a composition series of K. Writing \(G_i = \pi^{-1}(\mathbf{K}_i)\) for \(0 \leqslant i \leqslant p\) and \(G_i = H_{i-p}\) for \(p \leqslant i \leqslant n + p\) , a composition series \(\Sigma = (G_i)_{0 \leqslant i \leqslant n+p}\) of G is obtained. The sequence of quotients of \(\Sigma\) is obtained by juxtaposing the sequence of quotients of \(\Sigma''\) and the sequence of quotients of \(\Sigma'\) . If \(\Sigma'\) and \(\Sigma''\) are Jordan-Hölder series, \(\Sigma\) is a Jordan-Hölder series of G, by Proposition 9. If H or K admits composition series of arbitrary length, so does G. We have proved:

PROPOSITION 10. Let G be a group with operators and H a normal stable subgroup of G. The length of G is the sum of the lengths of H and G/H.

COROLLARY. Let G be a group with operators and \((\mathrm{G}_{i})_{0\leqslant i\leqslant n}\) a composition series of G. The length of G is the sum of the lengths of the \(G_{i}/G_{i+1}, 0 \leqslant i \leqslant n - 1\) .

If G and G' are isomorphic groups with operators and G admits a Jordan-Hölder series, so does G' and the Jordan-Hölder series of G and G' are equivalent. However, non-isomorphic groups can have equivalent Jordan-Hölder series; such is true for Z/4Z and (Z/2Z) × (Z/2Z), cf.~no. 10.

